\documentclass[11pt]{article}

\usepackage[margin=1in]{geometry}
\usepackage[T1]{fontenc}
\usepackage[utf8]{inputenc}
\usepackage{lmodern}
\usepackage{amsmath,amssymb,amsthm}
\usepackage{booktabs}
\usepackage{graphicx}
\usepackage{subcaption}
\usepackage{caption}
\usepackage{xcolor}
\usepackage{tcolorbox}
\usepackage{microtype}
\usepackage{url}
\usepackage[numbers,sort&compress]{natbib}
\usepackage[hidelinks]{hyperref}
\usepackage{cleveref}
\usepackage{algorithm}
\usepackage{algpseudocode}
\usepackage{fancyvrb}

\graphicspath{{figures/}}

% Without this, T1 lmtt renders "--" as an en-dash, silently corrupting every
% command-line flag in the paper.
\DisableLigatures[-]{encoding = T1, family = tt*}

\captionsetup{font=small,labelfont=bf}
\newcommand{\code}[1]{\texttt{\small #1}}
% Plain fancyvrb (breaklines/breakanywhere would need fvextra); the listings
% below are wrapped by hand with shell line continuations instead.
\DefineVerbatimEnvironment{shell}{Verbatim}%
  {fontsize=\small, xleftmargin=1em, samepage=false}
\newcommand{\GFM}{\textsc{GFM}}

% Gracefully degrade when a figure has not been generated yet.
\newcommand{\safeincludegraphics}[2][width=\linewidth]{%
  \IfFileExists{figures/#2.pdf}{\includegraphics[#1]{#2.pdf}}{%
    \IfFileExists{figures/#2.png}{\includegraphics[#1]{#2.png}}{%
      \fbox{\parbox[c][0.26\linewidth][c]{0.88\linewidth}{\centering\small\itshape
        figure \texttt{\detokenize{#2}} not generated yet}}}}}
\newcommand{\safeinput}[1]{\IfFileExists{#1}{\input{#1}}{\textit{(pending)}}}

\title{\bf Graph Flow Matching Beyond Faces and Bedrooms:\\
Porting \GFM{} to CIFAR-10, ImageNet and ImageNet-LT}

\author{%
  Nathan Sde\\
  \small Ben-Gurion University of the Negev\\
  \small \texttt{sde.nathan@gmail.com}
}
\date{\today}

\begin{document}
\maketitle

\begin{abstract}
Graph Flow Matching (\GFM{}) augments a latent flow-matching velocity field with
a graph-based correction term that couples every sample in a training batch to
its neighbours. The released implementation targets four
single-domain $256\times256$ datasets --- FFHQ, CelebA-HQ, LSUN Bedrooms and LSUN
Churches --- and hard-codes the corresponding latent resolution, dataset
loaders, and warm-start checkpoints. This report describes the engineering
required to run the \emph{same, unmodified} \GFM{} velocity network on three
substantially different datasets: CIFAR-10 (low resolution, ten balanced
classes), ImageNet-1k (high diversity, one thousand balanced classes) and
ImageNet-LT (the long-tailed subset, imbalance ratio $256{:}1$). We contribute a
resolution-agnostic latent pipeline, a labelled sharded-latent storage format
with chunk-aware batching, official/reconstructed ImageNet-LT split handling,
offline training-curve logging, and a single evaluation harness that reports
Fr\'echet Inception Distance (FID), Inception Score (IS), maximum mean
discrepancy (MMD), kernel Inception distance (KID) and a joint $t$-SNE embedding
of real and generated Inception features. \textbf{No component of the generative
model --- neither the reaction network, the graph adjacency generator, the
non-linear heat-diffusion correction, nor the flow-matching objective --- was
altered};  every change is confined to data handling, configuration and
measurement.
\end{abstract}

\safeinput{generated/status.tex}

\section{Introduction}
\label{sec:intro}

Flow matching~\citep{lipman2023flow,liu2023rectified} learns a time-dependent
velocity field $v_\theta(x,t)$ that transports a simple prior to the data
distribution along a prescribed probability path. Latent flow
matching~\citep{dao2023lfm} moves this construction into the latent space of a
pretrained variational autoencoder, which makes high-resolution generation
tractable. Graph Flow Matching~\citep{siddiqui2025gfm} observes that the
velocity field is normally predicted \emph{independently} for each sample, and
adds a correction term that lets samples within a batch exchange information
through a learned graph:
\begin{equation}
  \label{eq:gfm}
  v_\theta(x,t) \;=\; \underbrace{v^{\text{reac}}_{\theta_1}(x,t)}_{\text{per-sample reaction}}
  \;+\; \underbrace{v^{\text{corr}}_{\theta_2}(x,t)}_{\text{neighbour-aware correction}} .
\end{equation}
In the configuration used throughout this report the correction term is the
non-linear graph heat flow
\begin{equation}
  \label{eq:heat}
  v^{\text{corr}}_{\theta_2}(x,t)
  \;=\; -\,N_2\!\Big( G^{\!\top}\big[\sigma\big(G\,N_1(x,t)\big)\big],\; t \Big),
\end{equation}
where $G$ is the graph gradient operator induced by a batch adjacency matrix
$A(x,t) \in \mathbb{R}^{B\times B}$ produced by a learned attention module,
$(G z)_{ij} = \sqrt{A_{ij}}\,(z_i - z_j)$, and $N_1, N_2$ are small U-Nets. The
adjacency is built \emph{across the batch dimension}, so the correction is
resolution-agnostic by construction --- a property we exploit heavily below.

\paragraph{Why porting is non-trivial.}
Although \cref{eq:gfm,eq:heat} contain no dataset-specific assumption, the
released code base does. Concretely, the obstacles were:
\begin{enumerate}\itemsep2pt
  \item \textbf{Hard-coded latent geometry.} The latent grid $32\times32$ and the
  attention resolutions $\{32,16\}$ appear as literals in the model-construction
  code, in the sampler and in the FID probe.
  \item \textbf{Dataset-specific loaders only.} The four supported datasets each
  have a bespoke loader; there is no CIFAR-10 path, the ImageNet path is
  incomplete (its loader returns \code{(latent, label)} whereas the training loop
  destructures \code{(x0, x1)}, so the label would be consumed as the image), and
  there is no ImageNet-LT path at all.
  \item \textbf{Mandatory warm start.} Training defaults to
  \code{--use\_pretrained}, loading a public Latent-Flow-Matching checkpoint. No
  such checkpoint exists for our three datasets.
  \item \textbf{Mandatory in-training FID.} Checkpoints are written only inside
  the FID-evaluation branch, which itself requires pre-built \code{clean-fid}
  custom statistics. Without those, a run would train and save nothing.
  \item \textbf{No offline metrics.} Losses go only to Weights \& Biases, and
  neither IS, MMD nor any embedding visualisation exists in the repository.
\end{enumerate}

\paragraph{Contributions.}
\begin{itemize}\itemsep2pt
  \item A resolution-agnostic configuration path (\cref{sec:resolution}) that
  instantiates the \emph{unchanged} networks at whatever latent grid the data
  implies, together with a characterisation of the one latent-size limit that is
  genuinely baked into the released attention adjacency --- guarded up front
  rather than patched.
  \item A labelled, sharded latent format with a chunk-aware batch sampler
  (\cref{sec:storage}) that scales from CIFAR-10's $50$k images to ImageNet's
  $1.28$M.
  \item ImageNet-LT support with official-split download and a deterministic
  Pareto reconstruction fallback (\cref{sec:imagenetlt}).
  \item An evaluation harness (\cref{sec:protocol}) reporting FID, IS, MMD, KID
  and $t$-SNE from a single Inception feature pass, plus offline training-curve
  logging and plotting.
\end{itemize}

\section{Background}
\label{sec:background}

\subsection{Conditional flow matching}
Given data $x_1 \sim q$ and noise $x_0 \sim \mathcal{N}(0, I)$, the linear
(rectified) probability path interpolates
$z_t = (1-t)\,x_1 + t\,x_0$ under the reverse-time convention used by this code
base ($t=1$ is noise, $t=0$ is data), with target velocity
$u_t = x_0 - x_1$. Training minimises
\begin{equation}
  \label{eq:cfm}
  \mathcal{L}(\theta) \;=\;
  \frac{\mathbb{E}_{t,x_0,x_1}\big\|\, v_\theta(z_t,t) - u_t \,\big\|_2^2}
       {\mathbb{E}_{t,x_0,x_1}\big\|\, u_t \,\big\|_2^2},
\end{equation}
i.e.\ the mean-squared velocity error \emph{normalised by the target norm}. This
normalisation is what makes the loss values in \cref{sec:results} comparable
across datasets with different latent statistics; an unnormalised loss would be
dominated by whichever dataset has the larger latent variance.

\subsection{The graph correction term}
Within a batch of $B$ latents, an attention module with an ASPP trunk and a
sinusoidal time embedding emits a dense, non-negative adjacency
$A(x,t)\in\mathbb{R}^{B\times B}$. \cref{eq:heat} then performs one non-linear
diffusion step on the batch graph: $N_1$ maps each latent to a node signal, the
weighted graph gradient forms all $B^2$ pairwise differences, a pointwise
non-linearity $\sigma$ (ELU by default) is applied edge-wise, the weighted
divergence contracts back to $B$ node signals, and $N_2$ maps the result to a
velocity correction. Because every spatial dimension is flattened into the node
feature, the operator's cost is $O(B^2 \cdot C H W)$ and is \emph{independent of
how} $CHW$ is factorised --- the reason the same module works unchanged at
$4\times32\times32$ and $4\times16\times16$.

\section{Making \GFM{} dataset-portable}
\label{sec:port}

\begin{table}[t]
  \centering
  \small
  \begin{tabular}{p{0.30\linewidth} p{0.62\linewidth}}
    \toprule
    Component & Change \\
    \midrule
    Velocity network, adjacency generator, heat-flow correction, flow-matching
      loss, ODE solver
      & \textbf{None.} Used exactly as released. \\
    Model \emph{instantiation}
      & Latent grid and attention resolutions are read from the data instead of
        being literals. \\
    Data loading
      & New CIFAR-10 / ImageNet / ImageNet-LT datasets; new labelled sharded
        latent format; chunk-aware batch sampler. \\
    Training loop
      & Label-tolerant batch unpacking; FID-independent checkpointing; CSV
        logging of losses; from-scratch construction of the reaction network. \\
    Evaluation
      & New harness: FID, IS, MMD, KID, $t$-SNE, sample grids, curve plots. \\
    \bottomrule
  \end{tabular}
  \caption{Scope of the port. The generative model itself is untouched; every
  modification is configuration, data plumbing or measurement.}
  \label{tab:scope}
\end{table}

\subsection{Resolution-agnostic latent geometry}
\label{sec:resolution}

The Stable-Diffusion VAE (\code{stabilityai/sd-vae-ft-mse}) downsamples by a
factor of $8$ and emits four latent channels, so an input rendered at
$S_{\text{img}}$ pixels yields an $S = S_{\text{img}}/8$ latent grid. We expose
\code{--image\_size} and derive
\begin{equation*}
  S \;=\; S_{\text{img}}/8, \qquad
  \texttt{latent\_size} \leftarrow S, \qquad
  \texttt{attn\_resolutions} \leftarrow \{r \in R_0 : r \le S\},
\end{equation*}
clamping the configured attention resolutions so that a smaller latent grid
cannot silently produce an unusable U-Net. When training from pre-encoded
latents the values are instead read directly off the encoded dataset, so a
mismatch between the CLI defaults and the data on disk is impossible.

\paragraph{CIFAR-10 and the latent bottleneck.}
CIFAR-10 is natively $32\times32$. Encoding it at native resolution would give
$4\times4\times4$ latents --- too coarse for a patchifying transformer and far
outside the regime the VAE was trained for. We therefore bicubically upsample
CIFAR-10 to $S_{\text{img}}=256$, which restores exactly the $4\times32\times32$
latents the released configuration expects, and downsample generated images back
to $32\times32$ before measurement so that FID and IS follow the standard
CIFAR-10 protocol. This costs compute but keeps the model bit-for-bit identical.
A cheaper $S_{\text{img}}=128$ (a $16\times16$ latent grid) is available for the
resolution-agnostic adjacency modes, but not for the attention adjacency the
paper defaults to --- see the constraint below.

\paragraph{One genuine model-side constraint.}
Auditing the released code at other latent sizes surfaced a limit that
\emph{cannot} be lifted without editing the model, and which we therefore guard
rather than patch. The attention adjacency generator
(\code{AdjGenerator\_Attention\_ASPP\_Time}) applies a fixed
\code{AvgPool2d(4,4)} and then flattens a hard-coded $8\times8$ map, so its
projection matrix only lines up when $S = 32$. We verified that
$S \in \{8, 16, 64\}$ all fail with a shape error inside that projection,
whereas the \code{cosine} and \code{knn} adjacencies --- which operate directly
on the flattened latent --- run correctly at every $S$ we tested.
\code{train.py} now rejects the invalid combination up front with an actionable
message instead of crashing several layers deep. Consequently every experiment
reported here uses $S_{\text{img}} = 256 \Rightarrow S = 32$, exactly the
released geometry.

We additionally found that the \code{gaussian} adjacency is numerically
degenerate at these dimensionalities: with $\|x_i - x_j\|^2$ in the thousands,
$\exp\!\big(-\|x_i-x_j\|^2/2T^2\big)$ underflows to exactly zero, and the
$\sqrt{A_{ij}}$ inside the graph gradient then has an infinite derivative, so a
run NaNs on its first backward pass. This is a property of the released code
rather than of our port; \code{train.py} warns instead of silently proceeding.

\subsection{Labelled sharded latents and chunk-aware batching}
\label{sec:storage}

Pre-encoding the dataset once removes the VAE from the training loop entirely.
At $4\times32\times32$ in \code{fp16} a latent occupies $8$\,KiB, so CIFAR-10
needs $\approx0.4$\,GiB, ImageNet-LT $\approx0.9$\,GiB and full ImageNet
$\approx10$\,GiB --- the last of which does not fit in RAM. We write fixed-size
shards
\begin{center}\small
\code{chunk\_00000.pt} $=$ \{\code{latents}: $[n, C, S, S]$,
\code{labels}: $[n]$, \code{start\_idx}\}
\end{center}
alongside a \code{metadata.json} and a flat \code{labels.npy}. The reader keeps
an LRU cache of a few shards; a \emph{chunk-aware batch sampler} shuffles shard
order and shuffles within each shard, so a batch never straddles two shards and
the cache never thrashes. This generalises the strategy the repository already
used for LSUN Bedrooms, adding the class labels that the original format
discards.

\paragraph{Labels are carried, not consumed.}
All new loaders emit \code{(x0, x1, y)}. The training loop unpacks $2$- or
$3$-tuples and \emph{ignores} $y$: the flow objective in \cref{eq:cfm} stays
unconditional, exactly as in the released code. Carrying the label costs
nothing and is what makes the per-class $t$-SNE of \cref{sec:results} and any
future class-conditional study possible.

\subsection{ImageNet-LT}
\label{sec:imagenetlt}

ImageNet-LT~\citep{liu2019largescale} is a $115{,}846$-image subset of
ImageNet-1k whose per-class counts follow a Pareto distribution with
$\alpha = 6$, ranging from $1280$ images for the head class to $5$ for the tail
class (imbalance ratio $256{:}1$). The loader fetches the official
\code{ImageNet\_LT\_train.txt} file list from a list of mirrors and then
\emph{verifies} it against those published statistics before use; we confirmed
that the downloaded list is exactly $115{,}846$ images over $1000$ classes with
head $1280$ and tail $5$.

If every mirror is unreachable, the loader falls back to a \emph{deterministic
reconstruction} that reproduces the official count profile --- same $\alpha$,
same head and tail counts, same number of classes --- by subsampling the full
training split, with a shuffled class order so the head is not an artefact of
WordNet-ID ordering. The reconstruction yields $121{,}937$ images
($+5.3\%$ relative to the official split) with identical head and tail counts.
Which of the two was used is recorded in the encoded dataset's metadata and must
be reported, because the reconstruction matches the official split's
\emph{statistics} but not its exact file list.

\subsection{Training-loop robustness}
\label{sec:trainloop}

Three changes were needed before a run could complete at all:
\begin{enumerate}\itemsep2pt
  \item \textbf{From-scratch reaction networks.} With no public checkpoint for
  these datasets, \code{--use\_pretrained} is disabled automatically when the
  path does not exist, and from-scratch constructors were added for the ADM and
  ResNet reaction backbones (the DiT and PnP-U-Net ones already existed).
  \item \textbf{FID-independent checkpointing.} Checkpoints are now written
  every \code{--save\_every\_steps} steps regardless of whether the in-training
  FID probe is enabled, and \code{--cleanfid\_dataset\_name none} disables that
  probe cleanly.
  \item \textbf{Offline curves.} Per-step and per-epoch losses, the two term
  magnitudes, the learning rate and the FID probe are mirrored into CSV files in
  the run directory, so \cref{fig:loss} can be regenerated without a Weights \&
  Biases account.
\end{enumerate}

\section{Evaluation protocol}
\label{sec:protocol}

All metrics are computed from $N_{\text{gen}}$ generated images and a reference
set of real images that passed through \emph{exactly} the same
crop/resize pipeline as the training data --- a detail that dominates FID far
more than any modelling choice.

\paragraph{FID.} Fr\'echet Inception Distance~\citep{heusel2017gans} is computed
with \code{clean-fid}~\citep{parmar2022aliased}, which reproduces the reference
Inception graph and the correct resizing; the fallback path (torchvision
Inception-V3 pool features with the usual Fr\'echet formula) is reported
separately when \code{clean-fid} is unavailable, since the two are not
numerically interchangeable.

\paragraph{Inception Score.} IS~\citep{salimans2016improved} is
$\exp\big(\mathbb{E}_x\,\mathrm{KL}(p(y\mid x)\,\|\,p(y))\big)$, averaged over
$10$ splits, from the $1000$-way softmax of torchvision's Inception-V3. We also
report the IS of the \emph{real} reference set, which is the practical ceiling
for that dataset at that resolution and makes the generated score
interpretable --- an absolute IS is otherwise meaningless across resolutions.

\paragraph{MMD.} We use the unbiased U-statistic estimator with a
multi-bandwidth radial basis kernel on the $2048$-d Inception features:
\begin{equation}
  \label{eq:mmd}
  \widehat{\mathrm{MMD}}^2_u(X,Y) =
  \frac{1}{m(m-1)}\sum_{i\neq j} k(x_i,x_j)
  + \frac{1}{n(n-1)}\sum_{i\neq j} k(y_i,y_j)
  - \frac{2}{mn}\sum_{i,j} k(x_i,y_j),
\end{equation}
with $k(a,b) = \frac{1}{|\mathcal{S}|}\sum_{s\in\mathcal{S}}
\exp\!\big(-\|a-b\|^2/(2 s\sigma^2)\big)$, $\mathcal{S} = \{0.25, 0.5, 1, 2,
4\}$ and $\sigma^2$ set by the median heuristic on the pooled sample. Features
are standardised by the real set's per-dimension mean and standard deviation so
that the kernel is not dominated by the few Inception coordinates with the
largest raw scale. Because the full $N\times N$ Gram matrix is unnecessary at
these sample sizes, the estimate is averaged over $20$ random subsets of $2000$
images --- the same subset scheme KID conventionally uses --- and we report the
mean together with its across-subset standard deviation. We quote
$\widehat{\mathrm{MMD}} = \sqrt{\max(\widehat{\mathrm{MMD}}^2_u, 0)}$; the
clamp is needed because the unbiased estimator can go negative when the two
distributions are close.

\paragraph{KID.} For continuity with the generative-modelling literature we also
report KID~\citep{binkowski2018demystifying}, the same U-statistic with the cubic
polynomial kernel $k(a,b) = (a^\top b/d + 1)^3$ over $100$ subsets of $1000$
images. KID is unbiased and therefore far more reliable than FID at small
$N_{\text{gen}}$.

\paragraph{$t$-SNE.} Real and generated Inception features are standardised,
reduced to $50$ dimensions by PCA, and jointly embedded with
$t$-SNE~\citep{maaten2008visualizing} (perplexity $30$, PCA initialisation).
Embedding both sets \emph{jointly} is essential: separate embeddings are not
comparable, since $t$-SNE coordinates carry no absolute meaning. For CIFAR-10 we
additionally colour the real points by ground-truth class, which exposes whether
the unconditional model covers all ten modes or collapses onto a subset.

\paragraph{Sampling.} Samples are drawn by integrating the learned velocity
field with the repository's own \code{integrate\_ode} routine, unchanged, using
the same solver and step count the run was validated with (\cref{tab:config}).
The number of function evaluations is $4\times$ the step count for RK4.

\section{Experimental setup}
\label{sec:setup}

\begin{table}[t]
  \centering
  \small
  \safeinput{generated/config_table.tex}
  \caption{Training configuration per dataset, read back from each run's
  \code{run\_config.json}.}
  \label{tab:config}
\end{table}

All three datasets are pre-encoded once with \code{sd-vae-ft-mse} at
$S_{\text{img}} = 256$ and trained from scratch with the non-linear
heat-diffusion coupling and attention-based adjacency, i.e.\ the configuration
the original paper reports as its default. \cref{tab:config} lists the exact
settings recovered from each run.

\section{Results}
\label{sec:results}

\subsection{Training loss}

\Cref{fig:loss} shows the flow-matching objective of \cref{eq:cfm} over training
and \cref{tab:loss} summarises each run. Because the loss is divided by the
target-velocity norm it starts at exactly $1.0$ for an untrained network --- the
model predicting zero velocity --- so the curve reads directly as the fraction of
the target velocity the model still fails to explain, and is comparable across
datasets with very different latent statistics.

\Cref{fig:terms} separates the two additive terms of \cref{eq:gfm}. The graph
correction is expected to stay well below the per-sample reaction term: it is a
refinement of an already-reasonable velocity field, and a correction that grew
to dominate would indicate the reaction network had failed to learn. Tracking
both magnitudes is the cheapest available diagnostic that the graph branch is
contributing without destabilising training.


\begin{figure}[t]
  \centering
  \safeincludegraphics[width=\linewidth]{loss_curve}
  \caption{Flow-matching training loss (\cref{eq:cfm}). \emph{Left:} per-step
  loss, faint lines raw and solid lines exponentially smoothed. \emph{Right:}
  per-epoch mean. The loss is normalised by the target-velocity norm, so values
  are directly comparable across datasets.}
  \label{fig:loss}
\end{figure}

\begin{figure}[t]
  \centering
  \begin{subfigure}[t]{0.48\linewidth}
    \centering
    \safeincludegraphics[width=\linewidth]{loss_terms}
    \caption{Reaction vs.\ graph-correction magnitude.}
    \label{fig:terms}
  \end{subfigure}\hfill
  \begin{subfigure}[t]{0.48\linewidth}
    \centering
    \safeincludegraphics[width=\linewidth]{fid_curve}
    \caption{In-training FID probe.}
    \label{fig:fidcurve}
  \end{subfigure}
  \caption{Auxiliary training diagnostics. \cref{fig:terms} tracks
  $\max|\cdot|$ of the two additive terms of \cref{eq:gfm}; the correction term
  staying several orders of magnitude below the reaction term is the expected
  regime --- it is a \emph{correction}, not a replacement.}
  \label{fig:diag}
\end{figure}

\begin{table}[t]
  \centering
  \small
  \safeinput{generated/loss_table.tex}
  \caption{Training summary per dataset.}
  \label{tab:loss}
\end{table}

\subsection{Generative quality}

\Cref{tab:results} reports FID, IS, MMD and KID for each dataset, and
\cref{tab:reference} pairs the generated Inception Score with that of the real
reference set at the same resolution. The reference column matters: an absolute
IS is not interpretable across resolutions or preprocessing pipelines, so the
quantity to read is the gap between the two columns rather than either number
alone. The same logic applies to MMD, which is reported on features standardised
by the real set's statistics and is therefore on a comparable scale across rows.


\begin{table}[t]
  \centering
  \small
  \safeinput{generated/results_table.tex}
  \caption{Generative quality. FID and MMD lower is better; IS higher is better.
  MMD is $\sqrt{\max(\widehat{\mathrm{MMD}}^2_u,0)}$ on standardised Inception
  features (\cref{eq:mmd}). ``Res.'' is the resolution the metrics were computed
  at: $32$ for CIFAR-10 (standard protocol) and $256$ for the ImageNet
  variants.}
  \label{tab:results}
\end{table}

\begin{table}[t]
  \centering
  \small
  \safeinput{generated/reference_table.tex}
  \caption{Inception Score of the real reference set alongside the generated
  score. The real column is the practical ceiling at that resolution and makes
  the generated value interpretable.}
  \label{tab:reference}
\end{table}

\subsection{Feature-space geometry}

FID compresses the comparison to a single Gaussian moment-match, which hides
whether a model covers every mode or concentrates on a few. \Cref{fig:tsne}
embeds real and generated Inception features \emph{jointly}, so the two point
clouds live in one coordinate system and their overlap is meaningful: an
isolated cluster of generated points is a mode the model invents, and a region
of real points with no generated neighbours is a mode it misses.
\Cref{fig:tsneclass} makes this concrete for CIFAR-10 by colouring the real
points by ground-truth class --- since training is unconditional, whether the
generated cloud spreads across all ten class regions is exactly the question of
interest.


\begin{figure}[t]
  \centering
  \begin{subfigure}[t]{0.32\linewidth}
    \centering
    \safeincludegraphics[width=\linewidth]{tsne_cifar_10}
    \caption{CIFAR-10}
  \end{subfigure}\hfill
  \begin{subfigure}[t]{0.32\linewidth}
    \centering
    \safeincludegraphics[width=\linewidth]{tsne_imagenet}
    \caption{ImageNet}
  \end{subfigure}\hfill
  \begin{subfigure}[t]{0.32\linewidth}
    \centering
    \safeincludegraphics[width=\linewidth]{tsne_imagenet_lt}
    \caption{ImageNet-LT}
  \end{subfigure}
  \caption{Joint $t$-SNE embedding of Inception features for real (blue) and
  \GFM{}-generated (red) images. Overlap indicates that the generated
  distribution covers the real one; an isolated red cluster indicates a mode the
  model invents, and an uncovered blue region a mode it misses.}
  \label{fig:tsne}
\end{figure}

\begin{figure}[t]
  \centering
  \safeincludegraphics[width=0.62\linewidth]{tsne_cifar_10_byclass}
  \caption{CIFAR-10 $t$-SNE with real points coloured by ground-truth class and
  generated points in grey. Because the model is trained unconditionally, this
  is the direct test of whether it covers all ten modes.}
  \label{fig:tsneclass}
\end{figure}

\subsection{Qualitative samples}

\Cref{fig:samples} shows uncurated samples --- the first $64$ images the sampler
produced, not a selection --- decoded from the flow-matched latents with the same
frozen VAE used for encoding.


\begin{figure}[t]
  \centering
  \begin{subfigure}[t]{0.32\linewidth}
    \centering
    \safeincludegraphics[width=\linewidth]{samples_cifar_10}
    \caption{CIFAR-10}
  \end{subfigure}\hfill
  \begin{subfigure}[t]{0.32\linewidth}
    \centering
    \safeincludegraphics[width=\linewidth]{samples_imagenet}
    \caption{ImageNet}
  \end{subfigure}\hfill
  \begin{subfigure}[t]{0.32\linewidth}
    \centering
    \safeincludegraphics[width=\linewidth]{samples_imagenet_lt}
    \caption{ImageNet-LT}
  \end{subfigure}
  \caption{Uncurated samples, decoded from the flow-matched latents.}
  \label{fig:samples}
\end{figure}

\section{Discussion and limitations}
\label{sec:discussion}

\paragraph{The port is a portability study, not a tuning study.}
No hyper-parameter was optimised for the new datasets: learning rate, batch
size, network widths, solver and step count are inherited from the released
LSUN/CelebA configurations. Numbers in \cref{tab:results} should therefore be
read as \emph{what \GFM{} does out of the box} on these datasets, not as its
attainable ceiling. In particular, ImageNet and ImageNet-LT results are far more
sensitive to training budget than the single-domain datasets \GFM{} was
developed on.

\paragraph{Unconditional training on class-labelled data.}
ImageNet and ImageNet-LT are usually modelled class-conditionally, and the
state-of-the-art FIDs quoted for them assume classifier-free guidance. We train
unconditionally to keep the model unchanged, so our numbers are not comparable
to conditional ImageNet literature. The labels are carried through the pipeline,
so enabling the repository's existing conditional coupling is a configuration
change rather than new modelling work.

\paragraph{CIFAR-10 through a $256$-pixel VAE.}
Upsampling CIFAR-10 to $256$ and decoding back to $32$ passes every image
through a VAE round-trip it was not designed for. The reconstruction error of
that round-trip is a floor on the achievable FID that is independent of the flow
model. Reporting the reference-set IS (\cref{tab:reference}) partially exposes
this, but a VAE-reconstruction FID would isolate it exactly and is the natural
next diagnostic.

\paragraph{Batch-size coupling.}
The graph correction is computed over the training batch, so the adjacency
$A \in \mathbb{R}^{B\times B}$ --- and hence the function the model computes --- is
batch-size dependent. Comparisons across runs are only meaningful at equal $B$.
At sampling time the same coupling applies, which is why the evaluation harness
generates with a fixed batch size.

\paragraph{ImageNet-LT split provenance.}
If the official file list could not be fetched, the reported numbers use the
Pareto reconstruction of \cref{sec:imagenetlt}. It matches the official split's
count profile but not its file list, and results are consequently not directly
comparable to published ImageNet-LT numbers. The provenance string is stored in
the encoded dataset's metadata for exactly this reason.

\section{Reproducing this report}
\label{sec:repro}

The whole sequence is driven by \code{scripts/run\_pipeline.py}; the individual
stages are listed below so any one of them can be re-run in isolation. Swap
\code{cifar10} for \code{imagenet} or \code{imagenet-lt} and point
\code{--data\_root} at the ImageNet root that holds \code{train/} and
\code{val/}.

\paragraph{1. Pre-encode the dataset into sharded latents.}
\begin{shell}
python datasets/encode_new_datasets.py \
    --dataset cifar10 --data_root ./data/cifar10 \
    --out_root ./work/encoded/cifar10_256 \
    --image_size 256 --batch_size 64 --device cuda:0
\end{shell}

\paragraph{2. Build the reference set and the clean-fid statistics.}
\begin{shell}
python datasets/make_reference_set.py \
    --dataset cifar10 --data_root ./data/cifar10 \
    --out_dir ./work/reference/cifar10_32 \
    --source_image_size 256 --eval_image_size 32 \
    --num_images 50000 --cleanfid_name cifar10_gfm_32
\end{shell}

\paragraph{3. Train --- unchanged model, new configuration.}
\begin{shell}
python train.py --dataset cifar10 --use_pre_encoded \
    --encoded_dataset_path ./work/encoded/cifar10_256 \
    --image_size 256 --base_model dit --diffusion --adj_mode attention \
    --flow_model_type nonLinearHeatDiffusion2 \
    --model_savepath ./work/runs/cifar10/models \
    --image_savepath ./work/runs/cifar10/images \
    --train_batch_size 64 --flow_epochs 200 \
    --cleanfid_dataset_name none --device cuda:0
\end{shell}

\paragraph{4. Evaluate: FID, IS, MMD, KID, $t$-SNE, sample grid.}
\begin{shell}
python evaluation/evaluate_gfm.py \
    --checkpoint ./work/runs/cifar10/models/vel_net_best_fid.pt \
    --model_savepath ./work/runs/cifar10/models \
    --dataset cifar10 --run_name cifar10 \
    --reference_dir ./work/reference/cifar10_32 \
    --out_dir ./work/results/cifar10 \
    --num_samples 10000 --latent_size 32 --eval_image_size 32 \
    --cleanfid_dataset_name cifar10_gfm_32 --device cuda:0
\end{shell}

\paragraph{5. Plot the curves, fill the tables, typeset.}
\begin{shell}
python scripts/plot_training_curves.py \
    --run CIFAR-10=./work/runs/cifar10/models
python scripts/build_paper.py \
    --result CIFAR-10=./work/results/cifar10 \
    --run CIFAR-10=./work/runs/cifar10/models
cd paper && latexmk -pdf main.tex
\end{shell}

\bibliographystyle{plainnat}
\bibliography{refs}

\end{document}
